\documentclass[10pt]{article}

\usepackage[margin=0.75in]{geometry}
\usepackage{amsmath,amsthm,amssymb}
\usepackage{xcolor}
\usepackage{cancel}
\usepackage{graphicx}
\usepackage{changepage}
\usepackage{circuitikz}
\usepackage{pgfplots}
\usepackage{physics}
\usepackage{hyperref}
\usepackage{siunitx}
\usepackage{fontspec}
\usepackage{relsize}
\usepackage{subfig}
\usepackage{todonotes}
\usepackage{multicol, multirow, booktabs}
\usepackage[breakable]{tcolorbox}
\usepackage[inline]{enumitem}

\theoremstyle{definition}
\newtheorem{problem}{Problem}
\newtheorem{soln}{Solution}

\pgfplotsset{compat=newest}
\usetikzlibrary{lindenmayersystems}
\usetikzlibrary{arrows}
\usetikzlibrary{calc}
\usetikzlibrary{positioning, fit}
\usetikzlibrary{3d, perspective}

\definecolor{incolor}{HTML}{303F9F}
\definecolor{outcolor}{HTML}{D84315}
\definecolor{cellborder}{HTML}{CFCFCF}
\definecolor{cellbackground}{HTML}{F7F7F7}
\newcommand{\ui}{\hat{i}}
\newcommand{\uj}{\hat{j}}
\newcommand{\uk}{\hat{k}}
\newcommand{\ux}{\hat{x}}
\newcommand{\uy}{\hat{y}}
\newcommand{\uz}{\hat{z}}
\newcommand{\pr}[1]{{#1}^\prime}
\pgfdeclarelayer{background}  
\pgfsetlayers{background,main}
\AtBeginDocument{\RenewCommandCopy\qty\SI}
\newcommand{\justif}[2]{&{#1}&\text{#2}}

\makeatletter
\newcommand{\boxspacing}{\kern\kvtcb@left@rule\kern\kvtcb@boxsep}
\makeatother
\newcommand{\prompt}[4]{
    \ttfamily\llap{{\color{#2}[#3]:\hspace{3pt}#4}}\vspace{-\baselineskip}
}

\newcommand{\thevenin}[2]{
  \begin{center}
    \begin{circuitikz} \draw
      (0,0) -- (2,0) to[battery1, l_=$V_{Th}\eq#1$] (2,2) 
      to[resistor, l_=$R_{Th}\eq#2$] (0,2)
      ;
      \draw [o-] (-.07,2.079);
      \draw [o-] (-.07,0.079);
    \end{circuitikz}
  \end{center}
}

\newcommand{\norton}[2]{
  \begin{center}
    \begin{circuitikz} \draw
      (0,0) -- (3,0) to[american current source, l_=$I_{N}\eq#1$] (3,2) -- (0,2) (2,0)
      to[resistor, l=$R_{N}\eq#2$] (2,2)
      ;
      \draw [o-] (-.07,2.079);
      \draw [o-] (-.07,0.079);
    \end{circuitikz}
  \end{center}
}

\newcommand{\highlight}[1]{\colorbox{yellow}{$\displaystyle #1$}}

\newcommand{\ti}[1]{\widetilde{#1}}

\newfontface{\Kaufmann}{Kaufmann}
\DeclareTextFontCommand{\kf}{\Kaufmann}
\newcommand{\scriptr}{\fontsize{12pt}{12pt}\kf{r}}

\newfontface{\KaufmannB}{Kaufmann Bd BT}
\DeclareTextFontCommand{\kfb}{\KaufmannB}
\newcommand{\bscriptr}{\fontsize{12pt}{12pt}\kfb{r}}

\newcommand{\bv}[1]{\mathbf{#1}}

\title{Math 4360H: Assignment I}
\author{Jeremy Favro (0805980) \\ Trent University, Peterborough, ON, Canada}
\date{\today}

\begin{document}
\maketitle

% PROBLEM 1
\begin{problem}
\begin{enumerate}[label=(\alph*)]
  \item Let $a,b,c\in \mathbb{Z}$. Show that the linear Diophantine equation
        $$ax+by=c$$
        has a solution iff $\gcd(a,b)$ divides $c$.
  \item Find a solution to the linear Diophantine equation
        $$816x+2260y=44$$
        by applying the Euclidean Algorithm to $a=816$ and $b=2260$ to find their greatest common divisor.
  \item Suppose $(x_0,y_0)$ is a solution to $ax+by=c$, i.e. $ax_0+by_0=c$. Then
        $$ax+by=ax_0+by_0\iff a(x-x_0)=b(y_0-y).$$
        Let $d=\gcd(a,b)$, then
        $$\frac{a}{d}(x-x_0)=\frac{b}{d}(y_0-y).$$
        Since $\gcd(a/d,b/d)=1$ this implies $b/d$ divides $(x-x_0)$ and so
        $$x-x_0=k\left(\frac{b}{d}\right)\iff x=x_0+\left(\frac{b}{d}\right)$$
        for some $k\in \mathbb{Z}$. Subbing this into $a(x-x_0)=b(y_0-y)$ we get
        that $$ak\left(\frac{b}{d}\right)=b(y_0-y)\iff y=y_0-k\left(\frac{a}{d}\right)$$
        This shows that every solution has the form
        $$x=x_0+k\left(\frac{b}{d}\right),\quad y=y_0-k\left(\frac{a}{d}\right)$$
        Verify that this gives a solution for every $k\in\mathbb{Z}$, and therefore
        finds all solutions to $ax_0 + by_0 = c$.
        Find the complete set of solutions to each of the following Diophantine equations.
        \begin{enumerate}[label=(\roman*)]
          \item $816x+2260y=44$
          \item $42x+36y=25$
          \item $348x+228y=1044$
        \end{enumerate}
\end{enumerate}
\end{problem}
\begin{soln}
  \begin{enumerate}[label=(\alph*)]
    \item Suppose the equation has a solution. Then we have that
          $$ax_0+by_0=c.$$
          Since the left hand side of the above equation is divided by $\gcd(a,b)$ by the definition of a gcd, so must be the right hand side
          to preserve equality. Conversely, suppose that $\gcd(a,b)|c$. Then $c=\gcd(a,b)s$ for $s\in\mathbb{Z}$ and by the extended Euclidean algorithm we can write
          $$\gcd(a,b)=ax_0+by_0$$ with non-unique $(x_0,y_0)\in\mathbb{Z}\cross\mathbb{Z}$. Multiplying by $s$ then
          $$\gcd(a,b)s=c=ax_0 s+by_0 s$$ which gives a solution for $x=x_0s$, $y=y_0s$. Using part (b) as a bit of a hint we can then solve for these solutions in terms of $\gcd(a,b)$
          Since the $(x_0,y_0)$ are not unique we'll set up two equations,
          $$ax_0+by_0=c\quad\text{and}\quad ax_1+by_1=c$$
          and subtract them to obtain
          $$a(x_0-x_1)+b(y_0-y_1)=0\implies a(x_0-x_1)=b(y_1-y_0).$$
          which means $a$ divides $b(y_1-y_0)$ because $(x_0-x_1)\in\mathbb{Z}$, so we can pull out the common divisor
          to get that $a/\gcd(a,b)$ divides $y_1-y_0$ so by the division algorithm $$y_1=y_0+r\frac{a}{\gcd(a,b)}.$$
          Substituting this back into our equality gives
          $$a(x_0-x_1)=br\frac{a}{\gcd(a,b)}\implies x_1=x_0-r\frac{b}{\gcd(a,b)}$$
          (I didn't see this in part (c), whoops).

    \item \begin{align*}
            2260 & =2\cdot816+628                       \\
            \midrule
            816  & =1\cdot628+188                       \\
            \midrule
            628  & =3\cdot188+64                        \\
            \midrule
            188  & =2\cdot64+60                         \\
            \midrule
            64   & =1\cdot60+4                          \\
            \midrule
            60   & =15\cdot4+0\implies \gcd(816,2260)=4
          \end{align*}
          Since $4\mid 44$ the equation has solutions by part (a). Now we find a way to write $4=2260x_0+816y_0$, which we can do from the
          steps we did in the algorithm above,
          \begin{align*}
            4 & =64-60                                                      \\
              & =64-(188-2\cdot64)=3\cdot64-188                             \\
              & =3\cdot(628-3\cdot188)-188=3\cdot628-10\cdot188             \\
              & =3\cdot628-10\cdot(816-628)=-10\cdot816+13\cdot628          \\
              & =-10\cdot816+13\cdot(2260-2\cdot816)=13\cdot2260-36\cdot816
          \end{align*}
          Now we can mutliply through by $44/4=11$,
          $$143\cdot2260-396\cdot816=44.$$
          So our solution is $x=143$, $y=-396$.
    \item Substituting $$x=x_0+k\left(\frac{b}{d}\right),\quad y=y_0-k\left(\frac{a}{d}\right)$$
          into $ax+by$ we obtain
          $$a(x_0+k\left(\frac{b}{d}\right))+b(y_0-k\left(\frac{a}{d}\right))=ax_0+\cancel{ak\left(\frac{b}{d}\right)}+by_0-\cancel{bk\left(\frac{a}{d}\right)}=ax_0+by_0=c$$
          \begin{enumerate}[label=(\roman*)]
            \item We saw in (b) that $x_0=143$, $y_0=-396$ is a solution and that $\gcd(2260,816)=4$, so our solution set is
                  $$x=143+k\left(\frac{816}{4}\right)=143+204k,\quad y=-396-k\left(\frac{2260}{4}\right)=-396-565k,\, k\in\mathbb{Z}.$$
            \item Since $\gcd(42,36)=6\nmid 25$ we do not have integer solutions.
                  % \begin{align*}
                  %   42 & =36(1)+6 \\
                  %   36 & =6(6)
                  % \end{align*}
                  % so we have $\gcd(42,36)=6$ and we can write
                  % $$6=42-36$$
                  % so $$\frac{25}{6}(42-36=6)=\frac{}{}$$
            \item $348x+228y=1044$
                  \begin{align*}
                    348 & =228(1)+120 \\
                    \midrule
                    228 & =120(1)+108 \\
                    \midrule
                    120 & =108(1)+12  \\
                    \midrule
                    108 & =12(9)
                  \end{align*}
                  so $\gcd(348,228)=12\mid 1044$ so $\exists$ integer solutions. We can write
                  \begin{align*}
                    12 & =120-108                       \\
                       & =120-(228-120)=2(120)-228      \\
                       & =2(348-228)-228=-3(228)+2(348)
                  \end{align*}
                  which we can multiply through by $1044/12$ to obtain that
                  $$-261(228)+174(348)=1044$$
                  and hence that the full solution set is
                  $$x=174+k\left(\frac{228}{12}\right)=174+19k,\quad y=-261-k\left(\frac{348}{12}\right)=-261-29k,\, k\in\mathbb{Z}.$$
          \end{enumerate}
  \end{enumerate}
\end{soln}

% PROBLEM 2
\begin{problem}
Let $R$ be a ring that contains at least two elements. Suppose for each
nonzero $r\in R$, there is a unique element $\overline{r}\in R$  such that $r\overline{r}r=r$.
\begin{enumerate}[label=(\alph*)]
  \item Show that $R$ has no zero divisors.
  \item Show that $\overline{r}r\overline{r}=\overline{r}$.
  \item Show that $R$ has a multiplicative identity element.
  \item Show that $R$ is a division ring.
\end{enumerate}
\end{problem}
\begin{soln}
  \begin{enumerate}[label=(\alph*)]
    \item Let $r,s \in R$ such that $r\neq 0_R$ and $rs=0_R$. Then by definition we have an element $\overline{r}$ such that
          $r\overline{r}r=r$. Then,
          \begin{align*}
            r & = r\overline{r}r                                           \\
              & = r\overline{r}r + (rs)r \justif{\quad}{$rs=0$}            \\
              & = (r\overline{r} + (rs))r \justif{\quad}{Distributive law} \\
              & = r(\overline{r} + (s))r \justif{\quad}{Distributive law}
          \end{align*}
          by definition the element $\overline{r}$ which satisfies $r\overline{r}r=r$ is unique so $\overline{r}+s=\overline{r}\implies s=0$, hence
          there are no zero divisors.
    \item \begin{align*}
            \overline{r}r\overline{r}           & =\overline{r}r\overline{r}r\overline{r} \justif{\quad}{$r=r\overline{r}r$}   \\
            \implies r\overline{r}r\overline{r} & =r\overline{r}r\overline{r}r\overline{r}\justif{\quad}{Left multiply by $r$} \\
            \implies \overline{r}               & =\overline{r}r\overline{r}\justif{\quad}{$r\overline{r}r=r$}
          \end{align*}
    \item Here we rely on $R$ having left/right cancellation laws which it does by having no zero divisors. Let $x\in R$, $r\overline{r}r=r\implies xr\overline{r}r=xr$
          and we cancel the right $r$ to obtain that $x r\overline{r}=x$. We also have that $\overline{r}r\overline{r}=\overline{r}$ and, in similar fashion,
          $\overline{r}r\overline{r}=\overline{r}\implies \overline{r}r\overline{r}x=\overline{r}x\implies r\overline{r}x=x$ so
          $r\overline{r}x=xr\overline{r}=x$, hence $r\overline{r}=1_R$. This was troubling for a while until I realised that it's fine (I think), because,
          though $r\overline{r}$ is not unique, that is no problem as for any such combination we arrive at the same identity.
    \item If the previous proof holds, it is evident that all elements have inverses (and hence are units) as $r\overline{r}=1_R \implies r^{-1}=\overline{r}$. THis
          was all we needed to prove that $R$ is a division ring.
  \end{enumerate}
\end{soln}

% PROBLEM 3
\begin{problem}
Let $R$ be a set, with binary operations $+$ and $\cdot$ such that
\begin{enumerate}[label=(\arabic*)]
  \item $(R, +)$ is a group,
  \item $(R \setminus \left\{0_R\right\}, \cdot)$ is a group, (where $0_R$ is the additive identity element),
  \item For any $a, b, c \in R$, $a(b + c) = (ab) + (ac)$ and $(a + b)c =
          (ac) + (bc)$.
\end{enumerate}
Show that $(R, +, \cdot)$ is a division ring.
\end{problem}
\begin{soln}
  A ring is a structure with two operations, $+$ and $\cdot$, such that $(R,+)$ is an abelian group under addition, $\cdot$ is associative, and the distributive laws
  hold. A division ring is a ring where division is defined for every (nonzero) element, e.g. every (nonzero) element has a multiplicative inverse. Since $(R\setminus \left\{0_R\right\},\cdot)$ is a group
  we get associativity and multiplicative inverses for free, and we are given that $(R,+)$ is a group and that the distributive laws hold. Hence all we must show is that $(R,+)$ is abelian.
  Let $a,b\in R$. Then
  \begin{align*}
    (a+b)+(-1)\cdot(b+a) & =a+b-b-a \justif{\quad}{We have -1 from (2) and distributive law from (3)} \\
                         & =a-a\justif{\quad}{Additive inverses are defined from (1)}                 \\
                         & =0_R
  \end{align*}
  Hence $(R, +, \cdot)$ is a division ring.
\end{soln}
\end{document}